% Generated draft: requires mathematical and source review.
\documentclass[11pt]{article}
\usepackage[margin=25mm]{geometry}
\usepackage{fontspec}
\setmainfont{Latin Modern Roman}
\usepackage{amsmath,amssymb,mathtools,unicode-math}
\setmathfont{Latin Modern Math}
\usepackage{enumitem,needspace,xurl,hyperref}
\setlength{\parindent}{0pt}
\setlength{\parskip}{5pt}
\setlength{\emergencystretch}{3em}
\newcommand{\N}{\mathbb{N}}
\newcommand{\Z}{\mathbb{Z}}
\newcommand{\Q}{\mathbb{Q}}
\newcommand{\R}{\mathbb{R}}
\newcommand{\C}{\mathbb{C}}
\newcommand{\F}{\mathbb{F}}
\newcommand{\A}{\mathbb{A}}
\newcommand{\PP}{\mathbb{P}}
\newcommand{\E}{\mathbb{E}}
\newcommand{\eps}{\varepsilon}
\newcommand{\dd}{\,\mathrm{d}}
\begin{document}
\section*{UnsolvedMath ID: 5500019}
\textbf{Source problem code:} catalogue entry
\subsection{Definitions and mathematical statement}
Let $d\ge5$ and let $\mathcal S=\{S_1,\ldots,S_n\}$ be a family of surfaces in $\R^d$. Their arrangement subdivides the ambient space according to the surfaces. A vertical decomposition further subdivides arrangement cells by recursively introducing walls parallel to selected coordinate directions over lower-dimensional features, yielding elementary cells. The exact recursion, coordinate order, and treatment of degeneracies must be specified; the source supplies no definition of the decomposition convention.

The question asks for its complexity. Under the usual combinatorial interpretation, fix an admissible surface class $\mathcal C_d$ and a decomposition convention, let $v(\mathcal S)$ count its elementary cells or the specified total collection of faces, and seek the asymptotic order of
\[
 V_d(n)=\sup\{v(\mathcal S):\mathcal S\subset\mathcal C_d,\ |\mathcal S|=n\}.
\]
The source does not explicitly say whether it counts all cells of the full arrangement, the decomposition of a single cell or selected region, or construction time. It also does not define ``surface.'' The research summary discusses algebraic surfaces, but gives no uniform degree or description-complexity bound. Such a restriction is essential to a complexity estimate depending only on $n$ and $d$.

The intended problem is to obtain matching general upper and lower bounds in dimensions at least five, after those inputs are fixed. Whether $d$ is treated as fixed and what dependence is allowed in constants are not explicitly stated. No exact exponent is supplied for the general case.
\subsection{Short English statement}
Determine the worst-case complexity of vertical decompositions of $n$ surfaces in dimension $d\ge5$. The precise surface class and whether the full arrangement or a selected cell is decomposed must be fixed.

The supplied summary says the general case is open, with tight bounds only for specialized families. Its low-dimensional and hyperplane comparisons are qualitative, including the informal phrase ``$n^{d-2}$-ish,'' without precise hypotheses or a compatible specification of the counted object. They are not adopted as rigorous bounds for the full decomposition. The maintained TOPP entry is recorded as Open as accessed 2026-07-29; the OPEN-TRIAGE label is a supplied assessment, not independent verification.
\subsection{Sources}
\begin{itemize}
\item Local UnsolvedMath registry, permanent ID 5500019, source code catalogue entry, ``Vertical Decompositions in $\R^d$.'' Dataset curation and metadata: Ulam AI and UnsolvedMath Contributors, CC BY 4.0.
\item \emph{The Open Problems Project}, Computational Geometry, Problem 19, \url{https://topp.openproblem.net/p19}; source-TeX extraction, accessed 2026-07-29, maintained entry recorded as Open.
\item Public source specification: \url{https://github.com/edemaine/topp}. The supplied rights note says its repository licence expressly covers software; problem-text reuse remains NEEDS\_REVIEW. This qualification is separate from the dataset licence.
\item The research summary gives no primary references or precise statements for its low-dimensional and specialized-family comparisons.
\end{itemize}
\end{document}
